\documentclass[11pt]{article}
\usepackage[margin=1in]{geometry}
\usepackage{amsmath,amssymb,microtype,xcolor}
\usepackage[colorlinks=true,urlcolor=blue!55!black]{hyperref}

\title{Critical Benjamin--Ono--BBM Quasi-Invariance:\\A Literature Answer to arXiv:2103.04408}
\author{Automated mathematics run \texttt{fa\_banach\_001}, lane 7}
\date{17 August 2026}

\begin{document}
\maketitle

\begin{center}
\fbox{\parbox{0.9\textwidth}{\textbf{Status.}
Literature already answered: the core quasi-invariance statement at critical
dispersion is proved in arXiv:2107.03445.  The exact density formula and the
stronger all-finite-$p$ integrability conclusion remain open.}}
\end{center}

\section*{Source future-work statement}

Genovese, Luc\`a, and Tzvetkov study the fractional BBM equation
\[
 \partial_tu+\partial_t|D_x|^\beta u+\partial_xu+\partial_x(u^2)=0.
\]
Their Theorem 1.3 treats $\beta\in(1,2]$ and proves quasi-invariance of
exponentially cut-off Gaussian measures, all-finite-$p$ integrability of the
transported densities, and (under an additional regularity hypothesis) an
explicit density formula.  Immediately afterward, on PDF page 4, they write:
\begin{quote}
It is however known that (1.4) is globally well-posed for $\beta=1$ \dots\ and
we plan to study the extension of Theorem 1.3 to the case $\beta=1$ in a
future work.
\end{quote}
This is the sole extracted future-work signal in arXiv:2103.04408.

\section*{Decisive follow-up}

The same authors' later paper arXiv:2107.03445 is devoted precisely to the
$\beta=1$ periodic Benjamin--Ono--BBM equation.  Its introduction explicitly
states that it extends the previous $\beta>1$ achievements to the minimal
dispersion case $\beta=1$.  Theorem 1.1 (PDF page 2) proves that for every
$s>1$ and energy cutoff $R>0$, the cut-off Gaussian measure
$\widetilde\gamma_s$ is quasi-invariant under the flow and, for each time
$t$, its transported density belongs to $L^{p(t,R)}$ for some
$p(t,R)>1$.  Remark 1.2 further states that the energy cutoff can be removed
when only quasi-invariance of the reference Gaussian measure $\gamma_s$ is
required.

The supporting authors therefore knew they were answering the critical
dispersion continuation: the equation, parameter endpoint, authors, and
claimed extension are explicit matches.

\newpage
\section*{Scope limitation}

This is a complete literature answer to the \emph{core quasi-invariance}
part of the future-work target, not to every clause of the earlier Theorem
1.3.  The follow-up explicitly says that the minimal dispersion prevents its
method from identifying the Radon--Nikodym density; that question remains
open.  It also proves membership in some $L^p$, $p>1$, rather than in every
finite $L^p$.  Moreover, it works with an energy cutoff rather than the
earlier exponential $H^s$ cutoff, although Remark 1.2 yields the uncut
quasi-invariance conclusion.

\section*{Search evidence}

Run-index searches for the exact arXiv id, title, equation, and critical
parameter found no prior packet.  An exact-title arXiv query located
arXiv:2107.03445v2.  The supporting abstract, introduction, Theorem 1.1, and
Remark 1.2 were checked against the source's Theorem 1.3 and future-work
sentence.  The match is explicit rather than an agent-inferred reformulation.

\begin{thebibliography}{2}
\bibitem{source}
G. Genovese, R. Luc\`a, and N. Tzvetkov,
\emph{Transport of Gaussian measures with exponential cut-off for Hamiltonian PDEs},
arXiv:2103.04408v3; Ann. Inst. H. Poincar\'e C Anal. Non Lin\'eaire 39 (2022), 1465--1503.

\bibitem{answer}
G. Genovese, R. Luc\`a, and N. Tzvetkov,
\emph{Quasi-invariance of Gaussian measures for the periodic Benjamin--Ono--BBM equation},
arXiv:2107.03445v2; Stoch. Partial Differ. Equ. Anal. Comput. 10 (2022), 1482--1507.
\end{thebibliography}

\end{document}
